\section{Cómo convertirse en Agent Manager}

\subsection{La evolución de los equipos de desarrollo de software}

La forma de organización de los equipos de desarrollo de software está atravesando su cuarta transformación importante:

\begin{enumerate}
    \item \textbf{1940--1960s}: un solo desarrollador gestiona todo el proyecto
    \item \textbf{1970--1990s}: los equipos de software se vuelven la corriente principal, surge la división profesional del trabajo
    \item \textbf{2023--2025}: los desarrolladores se apoyan en sistemas de programación con IA para su trabajo
    \item \textbf{2025+}: un solo desarrollador gestiona el resultado del trabajo de varios Agentes de IA
\end{enumerate}

\begin{importantbox}{De gestionar personas a gestionar Agentes}
La ruta de evolución del desarrollo es: gestionar el propio resultado $\rightarrow$ gestionar el resultado de un equipo $\rightarrow$ gestionar el resultado de un equipo asistido por IA $\rightarrow$ \textbf{una sola persona gestiona el resultado de varios Agentes de IA}. Este último paso es la transformación que está ocurriendo en este momento.
\end{importantbox}

\subsection{Descomposición de las tareas de software}

Una tarea típica de desarrollo de software puede descomponerse en los siguientes pasos, donde el verde indica que la persona lidera y el azul indica que el Agente puede asumirlo:

\begin{enumerate}
    \item Aportar los requisitos de alto nivel \textbf{{[}persona{]}}
    \item Convertir los requisitos en un documento de diseño \textbf{{[}persona/Agente{]}}
    \item Implementar la solución según el documento de diseño \textbf{{[}Agente{]}}
    \item Añadir pruebas \textbf{{[}Agente{]}}
    \item Asegurar que el CI pase \textbf{{[}Agente{]}}
    \item Revisión de código \textbf{{[}Agente{]}}
    \item Actualizar la documentación \textbf{{[}Agente{]}}
\end{enumerate}

\begin{knowledgebox}{El valor central de la persona está en los pasos 1--2}
En el modelo de gestión de Agentes, el mayor valor del ingeniero humano se manifiesta en la \textbf{comprensión de los requisitos} y el \textbf{diseño de la arquitectura}. Estos dos pasos requieren contexto de negocio, criterio técnico y «gusto», que son justamente los eslabones más débiles de la IA actual.
\end{knowledgebox}

\subsection{Cuatro técnicas para dirigir a un Agente}

\subsubsection{Archivos de comportamiento del Agente}

Como CLAUDE.md, .cursorrules, AGENTS.md, entre otros, que aportan orientación de comportamiento a nivel de proyecto.

\subsubsection{Hooks}

Scripts deterministas que se ejecutan cuando se dispara un evento predefinido:
\begin{itemize}
    \item \texttt{PreToolUse}: antes de invocar una herramienta
    \item \texttt{PostToolUse}: después de invocar una herramienta
    \item \texttt{UserPromptSubmit}: cuando el usuario envía un prompt
    \item \texttt{PreCompact}: antes de comprimir el context
\end{itemize}

\subsubsection{Commands}

Guardar los prompts de uso frecuente como archivos, que el Agente puede ejecutar según se necesite. Casos de uso típicos: ejecutar pruebas, revisar código, hacer un commit de git y subirlo (push).

\subsubsection{Subagents}

Mecanismo de delegación en tiempo de ejecución, cuyos propósitos centrales son:
\begin{itemize}
    \item Crear «personalidades» de desarrollador independientes para distintos tipos de trabajo (frontend, backend, etc.)
    \item Separar con claridad el context de distintos flujos de trabajo
    \item Ofrecer un system prompt personalizado, un conjunto de herramientas propio y una context window independiente
\end{itemize}

\begin{importantbox}{El Subagent es el germen de que un Agente gestione a otros Agentes}
El modelo de subagent representa una tendencia importante: que un Agente gestione a otros Agentes. El Agente principal se encarga de descomponer y delegar las tareas, mientras que los subagents se concentran en la ejecución de un dominio específico; cada subagent cuenta con su propio context, lo que evita la sobrecarga de información.
\end{importantbox}

\subsection{Mejores prácticas}

\begin{enumerate}
    \item \textbf{Establecer barreras de protección estrictas}: las pruebas en el repositorio de código y las mejores prácticas de CI/CD son un respaldo (backstop) indispensable
    \item \textbf{Auditar cada modificación del Agente}: etiquetar (label) cada diff producido por un Agente
    \item \textbf{Usar modelos distintos para tareas distintas}: modelos rápidos para tareas simples, modelos potentes para tareas complejas
    \item \textbf{Las tareas complejas requieren más orientación previa}: no todas las tareas se prestan a una delegación completamente asíncrona
    \item \textbf{Hacer checkpoint (commit) con frecuencia}: confirmar el código de manera periódica para evitar perder avances cuando el Agente se desvíe del rumbo
\end{enumerate}

\begin{warningbox}{Preguntas abiertas}
\begin{itemize}
    \item ¿Cómo automatizar el 10--20\% inicial de la fase de investigación de cada tarea?
    \item ¿Cómo mantener una cola de tareas pendientes de ejecución? (Es relativamente fácil para una modificación puntual, pero sigue siendo un reto para un flujo de trabajo continuo)
\end{itemize}
\end{warningbox}

\subsection{Resumen de la sección}

El modelo de Agent Manager exige que el desarrollador se transforme de «la persona que escribe código» a «la persona que gestiona a los Agentes que escriben código». Las técnicas centrales incluyen los archivos de comportamiento, los hooks, los commands y los subagents. Los principios clave son: establecer barreras de protección, auditar las modificaciones, hacer coincidir la tarea con el modelo y hacer checkpoint con frecuencia.
